\documentclass[11pt]{article}

\usepackage[a4paper,margin=1.15in]{geometry}
\usepackage{amsmath,amssymb,amsthm,mathtools}
\usepackage{enumitem}
\usepackage{hyperref}
\usepackage{url}

\hypersetup{
  colorlinks=true,
  linkcolor=blue,
  citecolor=blue,
  urlcolor=blue
}

\numberwithin{equation}{section}

\newtheorem{theorem}{Theorem}[section]
\newtheorem{lemma}[theorem]{Lemma}
\newtheorem{proposition}[theorem]{Proposition}
\newtheorem{corollary}[theorem]{Corollary}
\newtheorem{claim}[theorem]{Claim}
\newtheorem{definition}[theorem]{Definition}
\newtheorem{remark}[theorem]{Remark}

\newcommand{\E}{\mathbb E}
\newcommand{\F}{\mathbb F}
\newcommand{\Z}{\mathbb Z}
\newcommand{\T}{\mathbb T}
\newcommand{\C}{\mathbb C}
\newcommand{\1}{\mathbf 1}
\newcommand{\wh}{\widehat}
\newcommand{\supp}{\operatorname{supp}}
\newcommand{\Span}{\operatorname{span}}
\newcommand{\CS}{\operatorname{CS}}
\newcommand{\calF}{\mathcal F}

\title{A Fourier--compactness argument for Erd\H{o}s Problem 42}
\author{X. L.}
\date{\today}

\begin{document}
\maketitle

\begin{abstract}
A proof of Erd\H{o}s Problem 42 is presented.  The argument reduces the
problem to a one-sided Fourier lemma for Cayley graphs over prime cyclic
finite fields.
The main technical point is a two-colour weighted degree-one
compactness/counting proposition.  Since the corresponding results in
Szegedy's paper are not stated in precisely this multicolour weighted form,
the required bridge is proved from the exact ingredients used there: the
ultraproduct Fourier factor, the degree-one generalized von Neumann
estimate, the structure of functions measurable with respect to the Fourier
sigma-algebra, and the dual embedding of degree-one limits.  The remaining
steps are the spectral passage from finite one-sided Fourier lower bounds
to Fourier positivity on the compact limit, the construction of a continuous
positive-definite representative, and a weighted positivity argument on
connected compact abelian groups.
\end{abstract}

\medskip
\noindent\textbf{2020 Mathematics Subject Classification.}
Primary 11B13; Secondary 11B30, 43A25, 05D10.

\medskip
\noindent\textbf{Keywords.}
Sidon sets; Erd\H{o}s problems; finite Fourier analysis; compact abelian groups;
ultraproducts; arithmetic regularity; Cauchy--Schwarz complexity.

\tableofcontents

\section{Introduction and statement of the theorem}

A finite set \(S\subset\Z\) is called a Sidon set if all pairwise sums
\[
        s+s',\qquad s,s'\in S,\quad s\le s',
\]
are distinct.  Equivalently, every nonzero integer has at most one
representation as an ordered difference \(s-s'\) with \(s,s'\in S\).

The problem considered here is the following question of Erd\H{o}s, from
\cite{Erdos1995}.  Bloom's Erd\H{o}s Problems entry is also cited for
the current problem statement and discussion record \cite{BloomErdos42}.

\begin{theorem}[Erd\H{o}s Problem 42]
\label{thm:erdos42}
For every \(M\ge 1\) there is \(N_0(M)\) such that, whenever
\(N\ge N_0(M)\) and \(A\subset\{1,\ldots,N\}\) is a Sidon set, there is a
Sidon set \(B\subset\{1,\ldots,N\}\) with
\[
        |B|=M
\]
and
\[
        (A-A)\cap(B-B)=\{0\}.
\]
\end{theorem}

\subsection*{Provenance and relation to the public discussion}
The proof below is written as a self-contained argument for every fixed
\(M\).  The public discussion attached to the Erd\H{o}s Problems page
contains relevant partial progress, reductions, and comments on the proposed
solution, and those contributions should be separated from the new part of
the present paper.  The problem page records the cases \(M=1,2,3\) as
partial progress in the comments and lists Zach Hunter and Daniil Sedov under
``Additional thanks'' \cite{BloomErdos42}; these partial results are not used
below.  In the discussion thread, Harjas posted the proposed solution that
motivated the present write-up \cite{HarjasProposedSolution}; qrdl formulated
a special case of the one-sided Fourier key lemma and the resulting reduction
to a random choice of \(B\) \cite{qrdlKeyLemma}; Thomas F. Bloom pointed out
the Cauchy--Schwarz complexity-one structure of the forms \(x_i-x_j\)
\cite{BloomComplexity}; Will Sawin described the regularity/compact-limit
interpretation of the key lemma \cite{SawinCompactLimit}; and Terence Tao
gave the continuous compact positive-definite analogue and explained how it
fits the compactness route \cite{TaoContinuousAnalogue}.  Tao also noted the
greedy reduction that removes the need to find a Sidon set directly
\cite{TaoGreedyReduction}; Appendix~\ref{app:sidon-subsets} gives the
elementary written proof of the extraction used here.  The main contribution
of this paper is the general fixed-\(M\) argument via
Lemma~\ref{lem:one-sided-cayley-fourier}, and especially the fully stated
multicolour compactness/counting bridge in
Proposition~\ref{prop:multicolour-compactness}.  The full discussion thread
is cited as \cite{Erdos42Discussion}; no unproved assertion from that thread
is used as a black box.


Finite cyclic groups are written additively.  For a prime \(p\), the
normalized Fourier transform on \(\F_p\) is
\[
        \widehat f(\gamma)
        =
        \frac1p\sum_{x\in\F_p}
        f(x)e^{-2\pi i\gamma x/p}.
\]
All expectations on finite groups are normalized.  On compact abelian
groups the convention is
\[
        \widehat f(\chi)=\int_G f(x)\overline{\chi(x)}\,dm_G(x).
\]

The central auxiliary result is the following one-sided Cayley--Fourier
lemma.

\begin{lemma}[One-sided Cayley--Fourier lemma]
\label{lem:one-sided-cayley-fourier}
Fix \(K\ge1\), \(0<\alpha<1\), and \(0<\theta<1\).  There are
\(\varepsilon>0\) and \(p_0\) such that the following holds for every
prime \(p\ge p_0\).

Let \(F\subset\F_p\) satisfy
\[
        0\in F,\qquad F=-F,\qquad |F|\le(1-\alpha)p,
\]
and
\[
        \widehat{\1_F}(\gamma)\ge-\varepsilon
        \qquad(\gamma\ne0).
\]
Set
\[
        I_p=[1,\lfloor\theta p\rfloor]\subset\F_p.
\]
Then
\[
        \E_{x_1,\ldots,x_K\in I_p}
        \prod_{1\le i<j\le K}
        \bigl(1-\1_F(x_i-x_j)\bigr)>0.
\]
Equivalently, there exist distinct \(x_1,\ldots,x_K\in I_p\) such that
\[
        x_i-x_j\notin F
        \qquad(i\ne j).
\]
\end{lemma}

\section{Reduction to the one-sided Cayley--Fourier lemma}

The first step is the reduction of Theorem~\ref{thm:erdos42} to
Lemma~\ref{lem:one-sided-cayley-fourier}.

\begin{lemma}[Fourier lower bound for Sidon difference sets]
\label{lem:sidon-fourier-lower}
Let \(A\subset\{1,\ldots,N\}\) be Sidon.  Let \(p>2N\) be prime and view
\(A\) as a subset of \(\F_p\).  Put
\[
        F=A-A\subset\F_p.
\]
Then
\[
        0\in F,\qquad F=-F,
\]
\[
        r_{A-A}=|A|\1_{\{0\}}+\1_{F\setminus\{0\}},
\]
\[
        |F|=|A|(|A|-1)+1,
\]
and, for every \(\gamma\in\widehat{\F_p}\),
\[
        \widehat{\1_F}(\gamma)
        \ge
        -\,\frac{|A|-1}{p}.
\]
Moreover
\[
        |A|\le 1+\sqrt{2N}
        \qquad\text{and}\qquad
        |F|\le 2N-1.
\]
\end{lemma}

\begin{proof}
Since \(A\subset[1,N]\), all integer differences \(a-a'\) lie in
\((-(N-1),N-1)\).  If two such differences are congruent modulo \(p\) and
\(p>2N\), then their difference is a multiple of \(p\) of absolute value
strictly less than \(p\), hence is zero.  Thus equality of differences
modulo \(p\) is exactly equality of differences in \(\Z\).

Let
\[
        r_{A-A}(x)=|\{(a,a')\in A^2:a-a'=x\}|.
\]
The diagonal pairs give \(r_{A-A}(0)=|A|\).  Since \(A\) is Sidon, every
nonzero difference has at most one ordered representation.  Hence
\[
        r_{A-A}=|A|\1_{\{0\}}+\1_{F\setminus\{0\}}.
\]
Equivalently,
\[
        \1_F=r_{A-A}-(|A|-1)\1_{\{0\}}.
\]
For the normalized Fourier transform,
\[
        \widehat{r_{A-A}}(\gamma)
        =
        p\,|\widehat{\1_A}(\gamma)|^2,
\]
and
\[
        \widehat{\1_{\{0\}}}(\gamma)=\frac1p.
\]
Therefore
\[
        \widehat{\1_F}(\gamma)
        =
        p|\widehat{\1_A}(\gamma)|^2-\frac{|A|-1}{p}
        \ge
        -\frac{|A|-1}{p}.
\]

The same difference representation shows that the \(|A|(|A|-1)\) ordered
nonzero differences are all distinct, and then
\[
        |F|=|A|(|A|-1)+1.
\]
Equivalently, considering only positive differences gives
\[
        \binom{|A|}{2}\le N-1.
\]
Thus \(|A|\le 1+\sqrt{2N}\), and
\[
        |F|=|A|(|A|-1)+1\le 2N-1.
\]
\end{proof}

\begin{proof}[Proof of Theorem~\ref{thm:erdos42} assuming
Lemma~\ref{lem:one-sided-cayley-fourier}]
Let \(R(M)\) be an integer such that every set of at least \(R(M)\)
integers contains a Sidon subset of size \(M\).  The existence of \(R(M)\)
is proved in Appendix~\ref{app:sidon-subsets}.

Apply Lemma~\ref{lem:one-sided-cayley-fourier} with
\[
        K=R(M),\qquad \alpha=\frac12,\qquad \theta=\frac18.
\]
Let \(\varepsilon>0\) and \(p_0\) be the corresponding constants.

Let \(N\) be sufficiently large in terms of \(M\), and let
\(A\subset[1,N]\) be Sidon.  By Bertrand's postulate, applied to \(4N\),
there is a prime
\[
        4N\le p\le 8N
\]
for all \(N\ge1\) after harmless adjustment of the endpoint convention.
In particular \(p>2N\).  View \(A\) as a subset of \(\F_p\), and set
\[
        F=A-A\subset\F_p.
\]
By Lemma~\ref{lem:sidon-fourier-lower},
\[
        |F|\le 2N-1<2N\le\frac p2.
\]
Also
\[
        \widehat{\1_F}(\gamma)
        \ge
        -\frac{|A|-1}{p}
        \ge
        -\frac{\sqrt{2N}}{4N}.
\]
For \(N\) sufficiently large this is at least \(-\varepsilon\).

Lemma~\ref{lem:one-sided-cayley-fourier} yields distinct elements
\[
        x_1,\ldots,x_K\in [1,\lfloor p/8\rfloor]\subset[1,N]
\]
such that
\[
        x_i-x_j\notin F
        \qquad(i\ne j)
\]
in \(\F_p\).  Let
\[
        X=\{x_1,\ldots,x_K\}.
\]
By the choice of \(K=R(M)\), \(X\) contains a Sidon subset \(B\) of size
\(M\).

It remains to pass from \(\F_p\) back to the interval.  If
\[
        b-b'=a-a'
\]
in \(\Z\), with \(b,b'\in B\) and \(a,a'\in A\), then the same equality
holds modulo \(p\).  Since \(B\subset X\), the construction gives
\[
        b-b'\notin A-A
\]
modulo \(p\) unless \(b=b'\).  Hence \(b=b'\), and then \(a-a'=0\), so
\(a=a'\).  Therefore
\[
        (A-A)\cap(B-B)=\{0\}.
\]
\end{proof}

\section{Proof of the one-sided Cayley--Fourier lemma by contradiction}

Assume Lemma~\ref{lem:one-sided-cayley-fourier} fails for some fixed
\[
        K\ge1,\qquad 0<\alpha<1,\qquad 0<\theta<1.
\]
Then there are primes \(p_n\to\infty\), numbers
\(\varepsilon_n\downarrow0\), and sets \(F_n\subset\F_{p_n}\) such that
\[
        0\in F_n,\qquad F_n=-F_n,\qquad |F_n|\le(1-\alpha)p_n,
\]
\[
        \widehat{\1_{F_n}}(\gamma)\ge-\varepsilon_n
        \qquad(\gamma\ne0),
\]
but, with
\[
        I_n=[1,\lfloor\theta p_n\rfloor]\subset\F_{p_n},
\]
one has
\[
        \E_{x_1,\ldots,x_K\in I_n}
        \prod_{i<j}
        \bigl(1-\1_{F_n}(x_i-x_j)\bigr)=0
\]
for every \(n\).

Set
\[
        u_n=\1_{I_n},
        \qquad
        h_n=\1_{F_n}.
\]
The multicolour compactness/counting proposition proved in
Section~\ref{sec:compactness} gives a connected compact abelian group
\(G\), functions \(u,h:G\to[0,1]\), and
\[
        \int_G u\,dm_G=\theta,
        \qquad
        \int_G h\,dm_G\le1-\alpha.
\]
The finite Fourier lower bound passes to the limit and makes \(h\)
Fourier-positive.  By the analytic lemmas in Section~\ref{sec:fourier-positive},
the function \(h\) may be replaced by a continuous positive-definite representative.
Then, for
\[
        d\nu=\frac{u\,dm_G}{\theta},
\]
the coloured counting statement gives
\[
\begin{aligned}
&\lim_{n\to\infty}
\E_{x_1,\ldots,x_K\in I_n}
\prod_{i<j}\bigl(1-h_n(x_i-x_j)\bigr)                                      \\
&\qquad =
\int_{G^K}
\prod_{i<j}\bigl(1-h(y_i-y_j)\bigr)\,d\nu^K(\mathbf y).
\end{aligned}
\]
The weighted positivity lemma in Section~\ref{sec:weighted-positivity}
shows that the last integral is strictly positive, contradicting the
vanishing of the finite averages.  Thus
Lemma~\ref{lem:one-sided-cayley-fourier} holds.

The remainder of the paper supplies the details used in this contradiction
argument.

\section{The two-colour degree-one compactness/counting proposition}
\label{sec:compactness}

This section is deliberately explicit about what is, and is not, taken
from Szegedy's paper \cite{SzegedyLimits}.  In Section 4, formula (1),
Szegedy defines the coloured density
\[
        t(L,F)
        =
        \E_{\mathbf x\in A^d}
        \prod_{\ell=1}^t f_\ell(L_\ell(\mathbf x)),
\]
where \(F=(f_1,\ldots,f_t)\) is a tuple of bounded functions.  Definition
4.1 defines true complexity through these coloured densities.  However,
Szegedy's Theorem 4 states continuity for the monochromatic map
\(f\mapsto t(L,f)\), and Proposition 6.2 is also stated for \(t(L,f)\).
Those two statements therefore do not, as written, directly give the
weighted two-colour proposition needed here.  The proof below supplies the
missing bridge.

The exact Szegedy ingredients used are the following.  Proposition 6.1
identifies linear characters of an ultraproduct as ultralimits of finite
linear characters.  Lemma 6.3 identifies the Fourier sigma-algebra
\(\calF(A)\) as the characteristic factor for the \(U^2\)-seminorm:
\(\|g\|_{U^2}=0\) if and only if \(g\) is orthogonal to
\(L^2(\calF(A))\).  Theorem 5 represents every bounded
\(\calF(A)\)-measurable function on a second-countable compact abelian
factor.  Corollary 7.2 gives the dual embedding for degree-one limits; in
the present prime-order situation this forces the limiting compact group
to be connected.  Proposition 6.2 is not used as a black box for the
coloured result; its proof motivates the factor reduction, which is spelled
out below.

\begin{definition}[Cauchy--Schwarz complexity at most one]
A finite system of integer linear forms
\[
        L=(L_1,\ldots,L_t)
\]
in variables \(x_1,\ldots,x_d\) has Cauchy--Schwarz complexity at most
\(1\) if, for every \(j\), the remaining forms can be partitioned into two
classes
\[
        \{L_i:i\ne j\}=\mathcal C_1\sqcup\mathcal C_2
\]
such that
\[
        L_j\notin\Span_{\mathbb Q}(\mathcal C_1)
        \quad\text{and}\quad
        L_j\notin\Span_{\mathbb Q}(\mathcal C_2).
\]
\end{definition}

\begin{lemma}[Generalized von Neumann estimate, complexity one]
\label{lem:gvn}
Let \(L=(L_1,\ldots,L_t)\) be a fixed system of integer linear forms of
Cauchy--Schwarz complexity at most \(1\).  Let \(A_n=\F_{p_n}\) with
\(p_n\to\infty\).  Then, after discarding finitely many \(n\), for every
\(j\) and every choice of functions \(f_{\ell,n}:A_n\to\C\) with
\(|f_{\ell,n}|\le1\),
\[
\left|
\E_{\mathbf x\in A_n^d}
\prod_{\ell=1}^t f_{\ell,n}(L_\ell(\mathbf x))
\right|
\le
\|f_{j,n}\|_{U^2(A_n)}.
\]
Consequently, if \(A=\prod_\omega A_n\) is a Szegedy ultraproduct and
\(f_\ell\in L^\infty(A)\) with \(|f_\ell|\le1\), then
\[
\left|
\int_{A^d}
\prod_{\ell=1}^t f_\ell(L_\ell(\mathbf x))\,d\mathbf x
\right|
\le
\|f_j\|_{U^2(A)}.
\]
\end{lemma}

\begin{proof}
For finite vector spaces this is the \(s=1\) case of the generalized von
Neumann theorem for systems of finite Cauchy--Schwarz complexity; see
Green--Tao \cite[Proposition 7.1]{GreenTaoPrimes}.  The only harmless
point in the present formulation is that the forms have integer
coefficients.  For a fixed system \(L\), the rational non-containment
conditions defining Cauchy--Schwarz complexity remain valid over
\(\F_{p_n}\) for all sufficiently large primes \(p_n\), because only
finitely many integer determinants are involved.

The quoted proposition is proved by applying Cauchy--Schwarz once for each
of the two classes in the defining partition for \(L_j\).  In this
complexity-one case the resulting expression is exactly bounded by the
\(U^2\)-norm of \(f_{j,n}\); all other functions remain bounded by \(1\).
Thus the displayed finite estimate follows.

For the ultraproduct statement, represent the bounded measurable
\(f_\ell\) up to null sets as ultralimits of bounded functions
\(f_{\ell,n}\) on \(A_n\), using the ultralimit measurability statement in
Szegedy's Section 5.  Apply the finite estimate and take the ultralimit.
The ultraproduct \(U^2\)-seminorm is, by definition, the ultralimit of the
finite \(U^2\)-norms for internal representatives and then extends by
approximation to bounded measurable functions.
\end{proof}

\begin{lemma}[Fourier-factor reduction for two colours]
\label{lem:two-colour-factor-reduction}
Let \(A=\prod_\omega\F_{p_n}\) be a Szegedy ultraproduct with
\(p_n\to\infty\), and let \(U,H\in L^\infty(A)\) satisfy
\(0\le U,H\le1\).  Let
\[
        U_0=\E(U\mid\calF(A)),
        \qquad
        H_0=\E(H\mid\calF(A)).
\]
For every Cauchy--Schwarz complexity-one system
\(L=(L_1,\ldots,L_t)\) and every colour assignment
\(C_\ell\in\{U,H\}\), with \(C_{\ell,0}=U_0\) or \(H_0\) according to the
same colour, one has
\[
        \int_{A^d}\prod_{\ell=1}^t C_\ell(L_\ell(\mathbf x))\,d\mathbf x
        =
        \int_{A^d}\prod_{\ell=1}^t C_{\ell,0}(L_\ell(\mathbf x))\,d\mathbf x.
\]
\end{lemma}

\begin{proof}
Write
\[
        C_\ell=C_{\ell,0}+R_\ell.
\]
By Szegedy's Lemma 6.3, each residual \(R_\ell\) is orthogonal to
\(L^2(\calF(A))\).  Hence
\[
        \|R_\ell\|_{U^2(A)}=0.
\]
Expanding the left-hand side multilinearly, the term with no residuals is
exactly the right-hand side.  Every other term contains at least one factor
\(R_j(L_j(\mathbf x))\).  All factors are bounded by \(1\), and
Lemma~\ref{lem:gvn} gives absolute value at most
\[
        \|R_j\|_{U^2(A)}=0.
\]
Thus all residual terms vanish.
\end{proof}

\begin{lemma}[Simultaneous representation on one compact factor]
\label{lem:simultaneous-factor}
Let \(A=\prod_\omega\F_{p_n}\) be as above, and let
\[
        U_0,H_0\in L^\infty(\calF(A)),
        \qquad 0\le U_0,H_0\le1.
\]
Then there are a second-countable compact abelian group \(G\), a
Szegedy-continuous surjective measure-preserving homomorphism
\[
        \pi:A\to G,
\]
and Borel functions \(u,h:G\to[0,1]\) such that
\[
        U_0=u\circ\pi,
        \qquad
        H_0=h\circ\pi
\]
almost everywhere.  Moreover \(G\) may be chosen connected.
\end{lemma}

\begin{proof}
Apply Szegedy's Theorem 5 to the single bounded
\(\calF(A)\)-measurable function
\[
        W_0=U_0+iH_0.
\]
It gives a second-countable compact abelian group \(G\), a
Szegedy-continuous surjective measure-preserving homomorphism
\(\pi:A\to G\), and \(w\in L^\infty(G)\) such that
\[
        W_0=w\circ\pi
\]
almost everywhere.  Set \(u=\operatorname{Re}w\) and
\(h=\operatorname{Im}w\), modifying on a Haar-null set so that
\(0\le u,h\le1\) everywhere.

It remains to record connectedness.  In the proof of Szegedy's Theorem 5,
\(G\) is generated by the countable family of characters occurring in the
Fourier expansion of \(W_0\), and the dual map
\[
        \widehat\pi:\widehat G\hookrightarrow\widehat A
\]
is injective.  By Szegedy's Proposition 6.1,
\[
        \widehat A=\prod_\omega\widehat{\F_{p_n}}
        \cong\prod_\omega\Z/p_n\Z.
\]
This ultraproduct has no nonzero torsion.  Indeed, if \(q\gamma=0\) and
\(\gamma=(\gamma_n)_\omega\), then \(q\gamma_n=0\) for \(\omega\)-almost
all \(n\); since \(p_n\to\infty\), \(q\) is invertible modulo \(p_n\) for
\(\omega\)-almost all \(n\), so \(\gamma_n=0\) for \(\omega\)-almost all
\(n\).  Hence \(\gamma=0\).  Thus \(\widehat G\) is torsion-free.  By
Pontryagin duality, a compact abelian group is connected if and only if
its dual group is torsion-free, so \(G\) is connected.
\end{proof}

\begin{proposition}[Multicolour degree-one compactness/counting]
\label{prop:multicolour-compactness}
Let \(p_n\to\infty\) be primes, let \(A_n=\F_{p_n}\), and let
\[
        u_n,h_n:A_n\to[0,1].
\]
After passing to a subsequence and fixing a non-principal ultrafilter
\(\omega\), there exist a single connected second-countable compact
abelian group \(G\), measurable functions
\[
        u,h:G\to[0,1],
\]
and a factor map
\[
        \pi:\prod_\omega A_n\to G
\]
with the following properties.

\begin{enumerate}[label=(\alph*)]
\item For every Cauchy--Schwarz complexity-one system
\(L=(L_1,\ldots,L_t)\) of integer linear forms and every colour assignment
\[
        c_{\ell,n}\in\{u_n,h_n\},
\]
with corresponding limit colour \(c_\ell\in\{u,h\}\), the finite coloured
densities converge and satisfy
\[
        \lim_{n\to\infty}
        \E_{\mathbf x\in A_n^d}
        \prod_{\ell=1}^t c_{\ell,n}(L_\ell(\mathbf x))
        =
        \int_{G^d}
        \prod_{\ell=1}^t c_\ell(L_\ell(\mathbf y))\,dm_G^d(\mathbf y).
\]

\item The same compact group \(G\) and the same factor map \(\pi\)
represent both colours: if
\[
        U=\lim_\omega u_n,
        \qquad
        H=\lim_\omega h_n,
\]
then
\[
        \E(U\mid\calF)=u\circ\pi,
        \qquad
        \E(H\mid\calF)=h\circ\pi
\]
almost everywhere on the ultraproduct, where \(\calF\) is Szegedy's
Fourier sigma-algebra.

\item The group \(G\) is connected; equivalently, \(\widehat G\) is
torsion-free.

\item For every \(\chi\in\widehat G\), there is a representing sequence of
characters \(\gamma_n\in\widehat{A_n}\) such that
\[
        \widehat h(\chi)
        =
        \lim_\omega \widehat{h_n}(\gamma_n).
\]
The analogous identity holds for \(u_n,u\).
\end{enumerate}
\end{proposition}

\begin{proof}
First pass to a subsequence along which every coloured density associated
with a Cauchy--Schwarz complexity-one integer system and colours from
\(\{u_n,h_n\}\) converges.  This diagonal selection is possible because
there are only countably many finite integer systems and finite colour
assignments.  It also makes the one-form densities \(\E u_n\) and
\(\E h_n\) converge.

Let
\[
        A=\prod_\omega A_n,
        \qquad
        U=\lim_\omega u_n,
        \qquad
        H=\lim_\omega h_n.
\]
Let \(\calF=\calF(A)\) be Szegedy's Fourier sigma-algebra, and put
\[
        U_0=\E(U\mid\calF),
        \qquad
        H_0=\E(H\mid\calF).
\]
Lemma~\ref{lem:two-colour-factor-reduction} shows that every degree-one
coloured density of \(U,H\) equals the corresponding density of
\(U_0,H_0\).  Lemma~\ref{lem:simultaneous-factor} gives one compact group
\(G\), one factor map \(\pi:A\to G\), and functions \(u,h:G\to[0,1]\) with
\[
        U_0=u\circ\pi,
        \qquad
        H_0=h\circ\pi.
\]
Since \(\pi\) is a measure-preserving homomorphism, \(\pi^d\) pushes the
ultraproduct probability measure on \(A^d\) to Haar measure on \(G^d\),
and
\[
        \pi(L_\ell(\mathbf x))
        =
        L_\ell(\pi(x_1),\ldots,\pi(x_d)).
\]
Thus the ultraproduct density equals the compact integral in part (a).
By the initial diagonal choice, the ordinary subsequential limit of the
finite densities equals the same ultralimit, proving (a).  Parts (b) and
(c) are exactly the content of Lemma~\ref{lem:simultaneous-factor}.

For (d), let \(\chi\in\widehat G\).  Then \(\chi\circ\pi\) is a linear
character of the ultraproduct.  By Szegedy's Proposition 6.1, there are
characters \(\gamma_n\in\widehat{A_n}\) such that
\[
        \chi\circ\pi=\lim_\omega \gamma_n.
\]
Using the Fourier convention stated in the introduction,
\[
\begin{aligned}
        \widehat h(\chi)
        &=\int_G h(y)\overline{\chi(y)}\,dm_G(y) \\
        &=\int_A H_0(x)\overline{\chi(\pi(x))}\,dx.
\end{aligned}
\]
Since \(H-H_0\) is orthogonal to \(L^2(\calF)\), and
\(\chi\circ\pi\) is \(\calF\)-measurable, the last integral equals
\[
        \int_A H(x)\overline{\chi(\pi(x))}\,dx.
\]
The ultraproduct integral of the internal function \(H\overline{\chi\circ\pi}\)
is
\[
        \lim_\omega \E_{x\in A_n}h_n(x)\overline{\gamma_n(x)}
        =
        \lim_\omega \widehat{h_n}(\gamma_n).
\]
The proof for \(u_n,u\) is identical.
\end{proof}

\begin{proposition}[Weighted compactness in the form used below]
\label{prop:weighted-compactness}
Let \(p_n\to\infty\) be primes and \(A_n=\F_{p_n}\).  Let
\[
        u_n=\1_{[1,\lfloor\theta p_n\rfloor]},
        \qquad
        h_n=\1_{F_n},
\]
where \(0<\theta<1\), \(F_n\subset A_n\), and
\[
        |F_n|\le(1-\alpha)p_n.
\]
After passing to a subsequence, Proposition~\ref{prop:multicolour-compactness}
gives a connected compact abelian group \(G\), Haar measure \(m_G\), a
factor map \(\pi\), and functions \(u,h:G\to[0,1]\) such that
\[
        \int_G u\,dm_G=\theta,
        \qquad
        \int_G h\,dm_G\le1-\alpha,
\]
and the coloured counting conclusion of
Proposition~\ref{prop:multicolour-compactness} holds for \(u_n,h_n\) and
\(u,h\).
\end{proposition}

\begin{proof}
Apply Proposition~\ref{prop:multicolour-compactness}.  Apply it first
to the one-form system
\[
        L_1(x)=x.
\]
For the colour \(u_n\), this gives
\[
        \int_G u\,dm_G
        =
        \lim_{n\to\infty}\E_{A_n}u_n
        =
        \lim_{n\to\infty}\frac{\lfloor\theta p_n\rfloor}{p_n}
        =
        \theta.
\]
For the colour \(h_n\), after passing to a further subsequence if necessary,
\[
        \frac{|F_n|}{p_n}\to\delta
\]
for some \(0\le\delta\le1-\alpha\).  Therefore
\[
        \int_G h\,dm_G
        =
        \lim_{n\to\infty}\E_{A_n}h_n
        =
        \delta
        \le1-\alpha.
\]
\end{proof}

\section{The coloured counting statement for the Cayley system}
\label{sec:cayley-counting}

\begin{lemma}[Hereditary property]
\label{lem:cs-hereditary}
Every subsystem of a Cauchy--Schwarz complexity-one system also has
Cauchy--Schwarz complexity at most \(1\).
\end{lemma}

\begin{proof}
Fix a form in the subsystem.  Take the two classes witnessing complexity
at most \(1\) for that form in the full system and intersect them with the
subsystem.  Passing to smaller classes can only shrink their rational
spans, so the required non-containment conditions remain true.
\end{proof}

\begin{lemma}[Complexity of the enlarged Cayley system]
\label{lem:cs-enlarged}
For every \(K\ge1\), the system
\[
        \mathcal L_K
        =
        \{x_1,\ldots,x_K\}
        \cup
        \{x_i-x_j:1\le i<j\le K\}
\]
has Cauchy--Schwarz complexity at most \(1\).
\end{lemma}

\begin{proof}
Fix \(L=x_a\).  Put
\[
        \mathcal C_1
        =
        \{x_a-x_j:j>a\}\cup\{x_i-x_a:i<a\},
\]
and
\[
        \mathcal C_2
        =
        \mathcal L_K\setminus(\{x_a\}\cup\mathcal C_1).
\]
Every form in \(\mathcal C_1\) is a difference, so every rational linear
combination of forms in \(\mathcal C_1\) has total coefficient sum zero.
The form \(x_a\) has total coefficient sum one, hence
\[
        x_a\notin\Span_{\mathbb Q}(\mathcal C_1).
\]
Every form in \(\mathcal C_2\) has coefficient zero in front of \(x_a\),
while \(x_a\) has coefficient one in front of \(x_a\).  Hence
\[
        x_a\notin\Span_{\mathbb Q}(\mathcal C_2).
\]

Now fix \(L=x_a-x_b\), with \(a<b\).  Let
\[
        \mathcal C_1
        =
        \{M\in\mathcal L_K\setminus\{x_a-x_b\}:
        \text{the coefficient of }x_b\text{ in }M\text{ is }0\},
\]
and
\[
        \mathcal C_2
        =
        \mathcal L_K\setminus(\{x_a-x_b\}\cup\mathcal C_1).
\]
Every rational linear combination of forms in \(\mathcal C_1\) has
coefficient zero in front of \(x_b\), whereas \(x_a-x_b\) has coefficient
\(-1\) in front of \(x_b\).  Thus
\[
        x_a-x_b\notin\Span_{\mathbb Q}(\mathcal C_1).
\]
The forms in \(\mathcal C_2\) are exactly those forms, different from
\(x_a-x_b\), with nonzero \(x_b\)-coefficient.  None of them involves
\(x_a\): the only form in \(\mathcal L_K\) involving both \(x_a\) and
\(x_b\) is \(x_a-x_b\), which has been removed.  Hence every rational
linear combination of forms in \(\mathcal C_2\) has coefficient zero in
front of \(x_a\), whereas \(x_a-x_b\) has coefficient one in front of
\(x_a\).  Therefore
\[
        x_a-x_b\notin\Span_{\mathbb Q}(\mathcal C_2).
\]
This proves the claim.
\end{proof}

\begin{lemma}[Coloured counting for the local Cayley average]
\label{lem:cayley-coloured-counting}
Let \(u_n,h_n,A_n,G,u,h\) be as in
Proposition~\ref{prop:weighted-compactness}.  Then
\[
\begin{aligned}
&\lim_{n\to\infty}
\E_{\mathbf x\in A_n^K}
        \prod_{r=1}^K u_n(x_r)
        \prod_{1\le i<j\le K}\bigl(1-h_n(x_i-x_j)\bigr) \\
&\quad =
\int_{G^K}
        \prod_{r=1}^K u(y_r)
        \prod_{1\le i<j\le K}\bigl(1-h(y_i-y_j)\bigr)
        \,dm_G^K(\mathbf y).
\end{aligned}
\]
\end{lemma}

\begin{proof}
Let \(P_K=\{(i,j):1\le i<j\le K\}\).  Expand
\[
        \prod_{(i,j)\in P_K}(1-h_n(x_i-x_j))
        =
        \sum_{S\subset P_K}(-1)^{|S|}
        \prod_{(i,j)\in S}h_n(x_i-x_j).
\]
For a fixed \(S\subset P_K\), the corresponding linear forms are
\[
        \{x_1,\ldots,x_K\}
        \cup
        \{x_i-x_j:(i,j)\in S\}.
\]
This is a subsystem of \(\mathcal L_K\), hence has Cauchy--Schwarz
complexity at most \(1\) by Lemmas~\ref{lem:cs-enlarged} and
\ref{lem:cs-hereditary}.  Proposition~\ref{prop:weighted-compactness}
therefore applies to each fixed \(S\), with colour \(u_n\) on the
coordinate forms and colour \(h_n\) on the selected difference forms:
\[
\begin{aligned}
&\lim_{n\to\infty}
\E_{\mathbf x\in A_n^K}
        \prod_{r=1}^K u_n(x_r)
        \prod_{(i,j)\in S}h_n(x_i-x_j) \\
&\quad =
\int_{G^K}
        \prod_{r=1}^K u(y_r)
        \prod_{(i,j)\in S}h(y_i-y_j)
        \,dm_G^K(\mathbf y).
\end{aligned}
\]
Summing the finitely many limits over \(S\subset P_K\) gives the displayed
identity.
\end{proof}

\section{Fourier positivity and continuous positive-definite representatives}
\label{sec:fourier-positive}

\begin{lemma}[Fourier positivity in the compact limit]
\label{lem:limit-fourier-positive}
Assume, in addition to the hypotheses of
Proposition~\ref{prop:weighted-compactness}, that
\[
        F_n=-F_n
\]
and
\[
        \widehat{h_n}(\gamma)\ge-\varepsilon_n
        \qquad(\gamma\ne0),
        \qquad
        \varepsilon_n\to0.
\]
Then the limiting function \(h\) satisfies
\[
        \widehat h(\chi)\ge0
        \qquad(\chi\in\widehat G).
\]
\end{lemma}

\begin{proof}
Let \(\chi\in\widehat G\).  By the spectral compatibility clause of
Proposition~\ref{prop:multicolour-compactness}, there are characters
\(\gamma_n\in\widehat{A_n}\) such that
\[
        \widehat h(\chi)=\lim_\omega \widehat{h_n}(\gamma_n).
\]
For \(\gamma=0\),
\[
        \widehat{h_n}(0)=\E h_n\ge0,
\]
and for \(\gamma\ne0\) the hypothesis gives
\[
        \widehat{h_n}(\gamma)\ge-\varepsilon_n.
\]
Hence the same lower bound \(-\varepsilon_n\) holds for all \(\gamma\).
Taking the ultralimit yields
\[
        \widehat h(\chi)
        \ge
        \lim_\omega(-\varepsilon_n)=0.
\]
Since \(F_n=-F_n\), all finite coefficients \(\widehat{h_n}(\gamma)\) are
real, and therefore the limiting coefficient is real and nonnegative.
\end{proof}

\begin{lemma}[Positive Fourier coefficients give a continuous representative]
\label{lem:linfty-continuous-pd}
Let \(G\) be a compact abelian group and let \(g\in L^\infty(G)\) satisfy
\[
        0\le g\le1
        \qquad\text{and}\qquad
        \widehat g(\chi)\ge0
        \quad(\chi\in\widehat G).
\]
Then
\[
        \sum_{\chi\in\widehat G}\widehat g(\chi)\le1,
\]
where the sum is the supremum over finite subsets of \(\widehat G\).
Consequently
\[
        \phi(x)=\sum_{\chi\in\widehat G}\widehat g(\chi)\chi(x)
\]
converges absolutely and uniformly, defines a continuous positive-definite
function \(\phi:G\to[0,1]\), and \(\phi=g\) almost everywhere.
\end{lemma}

\begin{proof}
Let \(S\subset\widehat G\) be finite and let \(\eta>0\).  By the
Bochner--Fejer kernel theorem for compact abelian groups, applied on the
finitely generated subgroup of \(\widehat G\) generated by \(S\), there is
a nonnegative trigonometric polynomial \(K\) on \(G\) such that
\[
        K\ge0,
        \qquad
        \int_G K\,dm_G=1,
\]
\[
        \widehat K(\rho)\ge0
        \quad(\rho\in\widehat G),
        \qquad
        \widehat K(\chi^{-1})\ge1-\eta
        \quad(\chi\in S).
\]
See Rudin \cite[Chapter 1, Sections 1.4--1.5]{Rudin} for the
Bochner--Fejer kernels and their approximation property.  Since
\(0\le g\le1\) and \(K\ge0\),
\[
        0\le\int_G gK\,dm_G\le1.
\]
Because \(K\) is a finite trigonometric polynomial,
\[
        \int_G gK\,dm_G
        =
        \sum_{\rho\in\widehat G}
        \widehat K(\rho)\widehat g(\rho^{-1}),
\]
with only finitely many nonzero summands.  All summands are nonnegative.
Therefore
\[
        (1-\eta)\sum_{\chi\in S}\widehat g(\chi)
        \le
        \int_G gK\,dm_G
        \le1.
\]
Letting \(\eta\downarrow0\) gives
\[
        \sum_{\chi\in S}\widehat g(\chi)\le1.
\]
Taking the supremum over finite \(S\) proves the asserted \(\ell^1\)
bound.

The Fourier series defining \(\phi\) is then absolutely and uniformly
convergent, so \(\phi\) is continuous.  For arbitrary
\(x_1,\ldots,x_m\in G\) and \(c_1,\ldots,c_m\in\C\),
\[
\begin{aligned}
        \sum_{i,j=1}^m c_i\overline{c_j}\phi(x_i-x_j)
        &=
        \sum_{\chi\in\widehat G}
        \widehat g(\chi)
        \left|
        \sum_{i=1}^m c_i\chi(x_i)
        \right|^2
        \ge0.
\end{aligned}
\]
Thus \(\phi\) is positive-definite.  Moreover
\[
        \widehat\phi(\chi)=\widehat g(\chi)
        \qquad(\chi\in\widehat G).
\]
All Fourier coefficients of \((g-\phi)dm_G\) vanish.  Since trigonometric
polynomials are uniformly dense in \(C(G)\), the measure
\((g-\phi)dm_G\) is zero, and \(g=\phi\) almost everywhere.
Finally, \(m_G\) has full support; continuity of \(\phi\) and the almost
everywhere bounds \(0\le g\le1\) imply \(0\le\phi\le1\) everywhere.
\end{proof}

\begin{lemma}[The one-set of a Fourier-positive function]
\label{lem:one-set}
Let \(g:G\to[0,1]\) be continuous and positive-definite, and suppose
\[
        g(x)=\sum_{\chi\in\widehat G} a_\chi\chi(x),
        \qquad
        a_\chi\ge0,
        \qquad
        \sum_\chi a_\chi\le1,
\]
with uniform convergence.  Put
\[
        H_g=\{x\in G:g(x)=1\},
        \qquad
        P=\{\chi\in\widehat G:a_\chi>0\}.
\]
If \(H_g\ne\varnothing\), then
\[
        H_g=\bigcap_{\chi\in P}\ker\chi.
\]
If \(H_g=\varnothing\), then the same conclusion is replaced by the
explicit alternative \(g(0)<1\).  In particular, whenever \(H_g\) is
nonempty, it is a closed subgroup of \(G\).  If additionally \(g\not\equiv1\),
then \(H_g\) is a proper closed subgroup.
\end{lemma}

\begin{proof}
If \(H_g=\varnothing\), then in particular \(g(0)<1\), since \(g(0)=1\)
would put \(0\in H_g\).  Now assume \(H_g\ne\varnothing\), and choose
\(x\in H_g\).  Then
\[
        1=|g(x)|
        \le
        \sum_\chi a_\chi
        =
        g(0)
        \le1.
\]
Equality holds throughout.  Equality in the triangle inequality implies
that every character with \(a_\chi>0\) has the same phase at \(x\); since
\(g(x)=1\) is real positive and \(\sum a_\chi=1\), that phase is \(1\).
Thus \(\chi(x)=1\) for every \(\chi\in P\), and
\[
        H_g\subset\bigcap_{\chi\in P}\ker\chi.
\]
Conversely, if \(y\in\bigcap_{\chi\in P}\ker\chi\), then
\[
        g(y)=\sum_\chi a_\chi=g(0)=1,
\]
so \(y\in H_g\).  This proves the equality and shows that \(H_g\) is a
closed subgroup when nonempty.  If \(H_g=G\), then \(g\equiv1\).  Hence,
if \(g\not\equiv1\), a nonempty \(H_g\) is proper.
\end{proof}

\begin{lemma}[Proper closed subgroups of connected compact groups]
\label{lem:proper-subgroup-null}
If \(G\) is a connected compact abelian group and \(H<G\) is a proper
closed subgroup, then \(m_G(H)=0\).
\end{lemma}

\begin{proof}
A closed subgroup of a compact group has positive Haar measure only if it
is open: this follows, for example, from the Steinhaus theorem, since
positive measure gives that \(HH^{-1}=H\) contains a neighbourhood of the
identity.  If \(H\) were open, then the quotient \(G/H\) would be a finite
discrete compact group.  The continuous image of a connected group is
connected, so \(G/H\) would be trivial, forcing \(H=G\).  Thus a proper
closed subgroup cannot have positive Haar measure.
\end{proof}

\begin{lemma}[Changing a compact representative on a null set]
\label{lem:null-representative}
Let \(u:G\to[0,1]\) be measurable with \(\theta=\int u\,dm_G>0\), and
let \(d\nu=u\,dm_G/\theta\).  If two Borel functions \(h,h':G\to[0,1]\)
agree Haar-almost everywhere, then for every \(K\),
\[
        \prod_{i<j}(1-h(y_i-y_j))
        =
        \prod_{i<j}(1-h'(y_i-y_j))
\]
for \(\nu^K\)-almost every \((y_1,\ldots,y_K)\).
\end{lemma}

\begin{proof}
Let \(E=\{x:h(x)\ne h'(x)\}\), so \(m_G(E)=0\).  Since \(\nu\ll m_G\),
every translate of \(E\) is \(\nu\)-null.  For fixed all variables except
\(y_i\), the set of \(y_i\) such that \(y_i-y_j\in E\) is a translate of
\(E\), hence \(\nu\)-null.  Fubini gives
\[
        \nu^K\{\mathbf y:y_i-y_j\in E\}=0
\]
for each pair \((i,j)\).  A finite union over pairs proves the claim.
\end{proof}

\section{Weighted positivity of the limiting integral}
\label{sec:weighted-positivity}

\begin{lemma}[Weighted positivity]
\label{lem:weighted-positivity}
Let \(G\) be a connected compact abelian group.  Let \(u:G\to[0,1]\) be
measurable and suppose
\[
        \theta:=\int_G u\,dm_G>0.
\]
Set
\[
        d\nu=\frac{u\,dm_G}{\theta}.
\]
Let \(h:G\to[0,1]\) be continuous and positive-definite, and suppose that
\[
        \int_G h\,dm_G\le1-\alpha
\]
for some \(\alpha>0\).  Then, for every \(K\ge1\),
\[
        \int_{G^K}
        \prod_{1\le i<j\le K}
        \bigl(1-h(y_i-y_j)\bigr)\,d\nu^K(\mathbf y)>0.
\]
\end{lemma}

\begin{proof}
For \(K=1\) the integral is \(1\).  Assume \(K\ge2\).
By Lemma~\ref{lem:linfty-continuous-pd}, applied to \(h\),
\[
        h(x)=\sum_{\chi\in\widehat G}a_\chi\chi(x),
        \qquad
        a_\chi\ge0,
        \qquad
        \sum_\chi a_\chi\le1,
\]
with uniform convergence.  Let
\[
        H_h=\{x\in G:h(x)=1\}.
\]
By Lemma~\ref{lem:one-set}, either \(H_h=\varnothing\), or \(H_h\) is a
closed subgroup.  Since \(\int h\,dm_G\le1-\alpha\), \(h\not\equiv1\), so
in the nonempty case \(H_h\) is a proper closed subgroup.  By
Lemma~\ref{lem:proper-subgroup-null},
\[
        m_G(H_h)=0
\]
whenever \(H_h\ne\varnothing\).  The same null conclusion is automatic if
\(H_h=\varnothing\).

Since \(\nu\ll m_G\), every translate of \(H_h\) is \(\nu\)-null.  Choose
\(y_1,\ldots,y_K\in\supp\nu\) inductively so that
\[
        y_i-y_j\notin H_h
        \qquad(i\ne j).
\]
At step \(r\), avoid the finite union
\[
        \bigcup_{i<r}(y_i+H_h),
\]
which has \(\nu\)-measure zero; because \(\nu(\supp\nu)=1\), this is
possible.  For the chosen tuple,
\[
        h(y_i-y_j)<1
        \qquad(i<j).
\]
By continuity of \(h\), there are neighbourhoods
\(V_1,\ldots,V_K\) of \(y_1,
\ldots,y_K\), and a number \(\eta>0\), such that
\[
        1-h(z_i-z_j)\ge\eta
        \qquad
        (z_i\in V_i,
        \ z_j\in V_j,
        \ i<j).
\]
Each \(V_i\) has positive \(\nu\)-measure because \(y_i\in\supp\nu\).
Therefore
\[
\begin{aligned}
        \int_{G^K}
        \prod_{i<j}(1-h(y_i-y_j))\,d\nu^K
        &\ge
        \eta^{\binom K2}\prod_{i=1}^K\nu(V_i)>0.
\end{aligned}
\]
\end{proof}

\section{Completion of the proof of the one-sided lemma}

The contradiction argument begun in Section 3 is now completed.

Assume Lemma~\ref{lem:one-sided-cayley-fourier} fails.  Then there are
\(p_n\to\infty\), \(F_n\subset\F_{p_n}\),
\(I_n=[1,\lfloor\theta p_n\rfloor]\),
\[
        u_n=\1_{I_n},
        \qquad
        h_n=\1_{F_n},
\]
and
\[
        \E_{x_1,\ldots,x_K\in I_n}
        \prod_{i<j}
        \bigl(1-h_n(x_i-x_j)\bigr)=0
\]
for all \(n\).

By Proposition~\ref{prop:weighted-compactness}, after passing to a
subsequence there are \(G,u,h\) with
\[
        \int_G u\,dm_G=\theta,
        \qquad
        \int_G h\,dm_G\le1-\alpha,
\]
and \(G\) connected.  By Lemmas~\ref{lem:limit-fourier-positive} and
\ref{lem:linfty-continuous-pd}, \(h\) may be replaced by a continuous
positive-definite representative with \(0\le h\le1\).  By
Lemma~\ref{lem:null-representative}, this replacement does not change the
weighted integrals below.

Rewrite the local average as a global weighted average:
\[
\begin{aligned}
&\E_{x_1,\ldots,x_K\in I_n}
        \prod_{i<j}
        \bigl(1-h_n(x_i-x_j)\bigr)                              \\
&\qquad =
\left(\frac{p_n}{|I_n|}\right)^K
\E_{\mathbf x\in\F_{p_n}^K}
        \prod_{r=1}^K u_n(x_r)
        \prod_{i<j}
        \bigl(1-h_n(x_i-x_j)\bigr).
\end{aligned}
\]
By Lemma~\ref{lem:cayley-coloured-counting},
\[
\begin{aligned}
&\lim_{n\to\infty}
\E_{\mathbf x\in\F_{p_n}^K}
        \prod_{r=1}^K u_n(x_r)
        \prod_{i<j}
        \bigl(1-h_n(x_i-x_j)\bigr)                              \\
&=\int_{G^K}
        \prod_{r=1}^K u(y_r)
        \prod_{i<j}
        \bigl(1-h(y_i-y_j)\bigr)\,dm_G^K.
\end{aligned}
\]
Since
\[
        \frac{|I_n|}{p_n}\to\theta,
\]
one obtains
\[
\begin{aligned}
&\lim_{n\to\infty}
\E_{x_1,\ldots,x_K\in I_n}
        \prod_{i<j}
        \bigl(1-h_n(x_i-x_j)\bigr)                              \\
&=
\theta^{-K}
\int_{G^K}
        \prod_{r=1}^K u(y_r)
        \prod_{i<j}
        \bigl(1-h(y_i-y_j)\bigr)\,dm_G^K                         \\
&=
\int_{G^K}
        \prod_{i<j}
        \bigl(1-h(y_i-y_j)\bigr)\,d\nu^K,
\end{aligned}
\]
where
\[
        d\nu=\frac{u\,dm_G}{\theta}.
\]
By Lemma~\ref{lem:weighted-positivity}, the last integral is strictly
positive.  This contradicts the assumption that all local averages are
zero.  Lemma~\ref{lem:one-sided-cayley-fourier} follows.

\section{Completion of Erd\H{o}s Problem 42}

Theorem~\ref{thm:erdos42} was reduced to
Lemma~\ref{lem:one-sided-cayley-fourier} in Section 2, and that lemma has
now been proved.  Hence, for every \(M\ge1\), choosing
\[
        K=R(M),\qquad \alpha=\frac12,
        \qquad \theta=\frac18
\]
and then \(N\) sufficiently large gives the desired Sidon set
\(B\subset[1,N]\) of size \(M\) with
\[
        (A-A)\cap(B-B)=\{0\}.
\]
This proves Theorem~\ref{thm:erdos42}.

\appendix

\section{Every sufficiently large finite set of integers contains a Sidon subset of size \texorpdfstring{\(M\)}{M}}
\label{app:sidon-subsets}

\begin{lemma}
\label{lem:large-set-sidon-subset}
For every \(M\ge1\) there is an integer \(R(M)\) such that every finite
set \(X\subset\Z\) with
\[
        |X|\ge R(M)
\]
contains a Sidon subset of size \(M\).
\end{lemma}

\begin{proof}
A greedy proof is given.  Suppose \(B\subset X\) is already Sidon and
\[
        |B|=t.
\]
It remains to estimate how many \(y\in X\setminus B\) cannot be added to \(B\) while
preserving the Sidon property.

The set \(B\cup\{y\}\) fails to be Sidon only if two sums involving
elements of \(B\cup\{y\}\) coincide and at least one of the two sums uses
\(y\).  Since \(B\) is Sidon, no equality among old sums is possible.
There are two types of forbidden equalities.

First, for \(b,b_1,b_2\in B\), one may have
\[
        y+b=b_1+b_2.
\]
For each ordered triple \((b,b_1,b_2)\in B^3\), this determines at most
one value
\[
        y=b_1+b_2-b.
\]
Thus this type forbids at most \(t^3\) values of \(y\).

Second, for \(b_1,b_2\in B\), one may have
\[
        2y=b_1+b_2.
\]
For each ordered pair \((b_1,b_2)\in B^2\), this determines at most one
integer \(y\).  Thus this type forbids at most \(t^2\) values.

The remaining possible equality
\[
        y+b=2y
\]
forces \(y=b\), and these values are already excluded by \(y\notin B\).
The equality \(y+b_1=y+b_2\) forces \(b_1=b_2\) and introduces no forbidden value.
Therefore, at stage \(t\), at most
\[
        t+t^2+t^3
\]
elements of \(X\) are unavailable, including the already chosen elements.
If
\[
        |X|>\sum_{t=0}^{M-1}(t+t^2+t^3),
\]
the greedy construction can be continued until it has chosen \(M\)
elements.  Thus one may take
\[
        R(M)=1+\sum_{t=0}^{M-1}(t+t^2+t^3).
\]
\end{proof}


\section*{Acknowledgements}
The author thanks Thomas F. Bloom for maintaining the Erd\H{o}s Problems
entry for Problem 42 and the contributors to its discussion thread for
public comments on partial cases, reductions, possible proof routes, and
attribution.  The author particularly acknowledges Harjas for posting the
proposed solution that motivated this write-up; qrdl for formulating the
one-sided Fourier key lemma and the reduction to a random avoiding set;
Thomas F. Bloom for the Cauchy--Schwarz complexity-one observation; Will
Sawin for comments on the regularity and compact-limit interpretation;
Terence Tao for the continuous compact positive-definite analogue of the key
lemma and for the greedy Sidon-extraction reduction; Kevin Barreto for
comments on attribution and on the compact Cayley-graph route; and Zach
Hunter and Daniil Sedov, who are listed on the problem page under
``Additional thanks''.  The author is solely responsible for the proofs and
for the verification of the compactness/counting argument in
Section~\ref{sec:compactness}.

\begin{thebibliography}{99}

\bibitem{Erdos1995}
P. Erd\H{o}s,
\emph{Some of my favourite problems in number theory, combinatorics, and geometry},
Resenhas do Instituto de Matem\'atica e Estat\'istica da Universidade de S\~ao Paulo
\textbf{2} (1995), no. 2, 165--186.

\bibitem{BloomErdos42}
T. F. Bloom,
\emph{Erd\H{o}s Problem \#42},
\url{https://www.erdosproblems.com/42},
accessed 2 May 2026.

\bibitem{Erdos42Discussion}
Erd\H{o}s Problems,
\emph{Discussion thread for Erd\H{o}s Problem \#42},
\url{https://www.erdosproblems.com/forum/thread/42},
accessed 2 May 2026.

\bibitem{HarjasProposedSolution}
Harjas,
\emph{Proposed solution comment on Erd\H{o}s Problem \#42},
Erd\H{o}s Problems discussion thread, 27 April 2026,
\url{https://www.erdosproblems.com/forum/thread/42},
accessed 2 May 2026.

\bibitem{qrdlKeyLemma}
qrdl,
\emph{Comment formulating a one-sided Fourier key lemma for Erd\H{o}s Problem \#42},
Erd\H{o}s Problems discussion thread, 29 April 2026,
\url{https://www.erdosproblems.com/forum/thread/42},
accessed 2 May 2026.

\bibitem{BloomComplexity}
T. F. Bloom,
\emph{Comment on Cauchy--Schwarz complexity one for the forms \(x_i-x_j\)},
Erd\H{o}s Problems discussion thread for Problem \#42, 29 April 2026,
\url{https://www.erdosproblems.com/forum/thread/42},
accessed 2 May 2026.

\bibitem{SawinCompactLimit}
W. Sawin,
\emph{Comment on the regularity and compact-limit interpretation of the key lemma},
Erd\H{o}s Problems discussion thread for Problem \#42, 29 April 2026,
\url{https://www.erdosproblems.com/forum/thread/42},
accessed 2 May 2026.

\bibitem{TaoContinuousAnalogue}
T. Tao,
\emph{Comment giving the continuous compact positive-definite analogue of the key lemma},
Erd\H{o}s Problems discussion thread for Problem \#42, 29 April 2026,
\url{https://www.erdosproblems.com/forum/thread/42},
accessed 2 May 2026.

\bibitem{TaoGreedyReduction}
T. Tao,
\emph{Comment on reductions for Erd\H{o}s Problem \#42, including greedy Sidon extraction},
Erd\H{o}s Problems discussion thread for Problem \#42, 5 December 2025,
\url{https://www.erdosproblems.com/forum/thread/42},
accessed 2 May 2026.

\bibitem{BarretoCompactCayley}
K. Barreto,
\emph{Comment on attribution and the compact Cayley-graph route for Erd\H{o}s Problem \#42},
Erd\H{o}s Problems discussion thread for Problem \#42, 1 May 2026,
\url{https://www.erdosproblems.com/forum/thread/42},
accessed 2 May 2026.

\bibitem{GreenTaoPrimes}
B. Green and T. Tao,
\emph{Linear equations in primes},
Annals of Mathematics \textbf{171} (2010), no. 3, 1753--1850.

\bibitem{Rudin}
W. Rudin,
\emph{Fourier Analysis on Groups},
Interscience Tracts in Pure and Applied Mathematics, No. 12,
Interscience Publishers, 1962.

\bibitem{SzegedyLimits}
B. Szegedy,
\emph{Limits of functions on groups},
arXiv:1502.07861, 2015.

\end{thebibliography}

\end{document}
